\chapter{คู่มือการติดตั้งและการพัฒนาสภาพแวดล้อม OpenXR}
\label{chap:appendix_a}

\section{การติดตั้ง OpenXR Software Development Kit (SDK)}

การพัฒนาแอปพลิเคชันบนมาตรฐาน OpenXR ต้องอาศัยการติดตั้งตัวโหลด (OpenXR Loader) และเฮดเดอร์ไฟล์มาตรฐานจาก Khronos Group

\subsection{การติดตั้งบนระบบปฏิบัติการ Linux (Ubuntu / Debian)}
บนระบบปฏิบัติการ Linux สามารถติดตั้งผ่านตัวจัดการแพ็กเกจได้โดยตรง
\begin{codebox}[คำสั่งติดตั้ง OpenXR SDK และ Monado Runtime บน Ubuntu]
sudo apt-get update
sudo apt-get install -y libopenxr-dev openxr-layer-apidump
sudo apt-get install -y monado-cli monado-gui
\end{codebox}

\subsection{การติดตั้งบนระบบปฏิบัติการ Windows}
สำหรับระบบปฏิบัติการ Windows สามารถดาวน์โหลด OpenXR SDK ผ่าน vcpkg หรือติดตั้งผ่าน Windows Mixed Reality Portal และ Meta Quest Link ซึ่งจะมี OpenXR Loader และ Runtime มาพร้อมกับระบบปฏิบัติการ

\section{การกำหนดค่าสภาพแวดล้อม Python และ PyOpenXR}

ในการทำปฏิบัติการตามตำราเล่มนี้ แนะนำให้สร้างสภาพแวดล้อมเสมือนของ Python (Virtual Environment) เพื่อป้องกันความขัดแย้งของแพ็กเกจ

\begin{codebox}[การสร้างสภาพแวดล้อมและติดตั้งไลบรารีที่จำเป็น]
# 1. สร้างสภาพแวดล้อมเสมือน
python3 -m venv openxr_stem_env
source openxr_stem_env/bin/activate

# 2. ปรับปรุง pip และติดตั้งแพ็กเกจหลัก
pip install --upgrade pip
pip install pyopenxr numpy scipy matplotlib
pip install PyOpenGL PyOpenGL_accelerate glfw
pip install mediapipe opencv-python
\end{codebox}

\section{การตรวจสอบความพร้อมของรันไทม์ด้วยคำสั่ง CLI}

หลังจากติดตั้งเสร็จสมบูรณ์ สามารถใช้คำสั่ง \texttt{hello\_xr} หรือคำสั่งทดสอบเพื่อตรวจสอบว่าระบบสามารถตรวจพบ OpenXR Runtime ได้อย่างถูกต้อง
\begin{codebox}[การทดสอบรันไทม์ผ่านคำสั่ง hello\_xr]
hello_xr -g OpenGL
\end{codebox}

หากระบบแสดงหน้าต่างสามมิติพร้อมคิวบ์หมุน แสดงว่าการติดตั้ง OpenXR Loader และ Graphic Swapchain ทำงานได้อย่างสมบูรณ์และพร้อมสำหรับการพัฒนาห้องปฏิบัติการเสมือนจริงในบทเรียนต่าง ๆ
\clearpage
